\documentclass[11pt]{article}
\usepackage[utf8]{inputenc}
\usepackage[T1]{fontenc}
\usepackage{amsmath,amssymb,mathtools}
\usepackage[margin=1in]{geometry}
\usepackage{hyperref}
\usepackage{array}
\usepackage{parskip}
\setlength{\parindent}{0pt}
\hypersetup{colorlinks=true,linkcolor=blue,urlcolor=blue}
\title{Lecture 22 — Sampling}
\author{}
\date{}
\begin{document}
\maketitle

\textbf{Course:} CEC 315 — Signals and Systems (Spring 2026)
\textbf{Instructor:} Rogelio Gracia Otalvaro
\textbf{Source PDF:} \texttt{all\_lectures/cec315-lctr22-sampling.pdf}
\textbf{Exam coverage:} Exam 3

\textbf{Lctr 22: Sampling}
Spring 2026
Focus: Chapter 7 (Pages 514–555)

\section*{1 The Big Picture: From Continuous Time to Discrete Time}

\subsection*{Why This Matters}

Nearly every modern signal-processing system converts continuous-time signals into discrete-time signals at some point. Your phone digitizes your voice before compressing and transmitting it. A digital camera samples a spatially continuous scene onto a finite grid of pixels. An MRI scanner collects discrete frequency-domain measurements of a continuous spatial distribution.

The central question of this lecture is: \textbf{Under what conditions can we recover a continuous-time signal perfectly from a set of equally spaced samples?} The answer is the \textit{sampling theorem}, one of the most important results in all of signals and systems.

\textbf{Roadmap for this lecture:}

\begin{enumerate}
  \item Introduce impulse-train sampling and derive its frequency-domain effect.
  \item State and interpret the sampling theorem.
  \item Show how to reconstruct a signal from its samples (band-limited interpolation).
  \item Explore what happens when we violate the sampling theorem (aliasing).
  \item Connect sampling to discrete-time processing of continuous-time signals.
\end{enumerate}

\section*{2 Impulse-Train Sampling}

\subsection*{2.1 The Sampling Model}

To study sampling mathematically, we model the process as multiplying $x(t)$ by a periodic impulse train $p(t)$ with period $T$:

\[x_p(t) \;=\; x(t)\,p(t), \qquad p(t) \;=\; \sum_{n=-\infty}^{+\infty} \delta(t-nT). \tag{1}\]

The period $T$ is the \textbf{sampling period} (time between samples), and the fundamental frequency $\omega_s = 2\pi/T$ is the \textbf{sampling frequency} (how many radians per second we sample at).

Because multiplying by an impulse simply captures the signal value at that instant ($x(t)\,\delta(t-nT) = x(nT)\,\delta(t-nT)$), the sampled signal is an impulse train whose amplitudes are the sample values:

\[x_p(t) \;=\; \sum_{n=-\infty}^{+\infty} x(nT)\,\delta(t-nT). \tag{2}\]

Think of it this way: we "grab" the value of $x(t)$ at every multiple of $T$ and store it as the height of an impulse at that location.

\textit{[Figure: three-panel diagram showing original signal $x(t)$ (smooth curve) times impulse train $p(t)$ (uniformly spaced impulses of period $T$) equals sampled signal $x_p(t)$ (impulses whose heights match $x$ at the sample times).]}

\subsection*{2.2 Frequency-Domain Effect of Sampling}

From the multiplication property of the Fourier transform, multiplying in time corresponds to convolution in frequency (scaled by $1/2\pi$). The Fourier transform of the impulse train is itself an impulse train:

\[P(j\omega) \;=\; \frac{2\pi}{T}\sum_{k=-\infty}^{+\infty}\delta(\omega-k\omega_s). \tag{3}\]

Convolving $X(j\omega)$ with this impulse train produces shifted replicas:

\[\boxed{\,X_p(j\omega) \;=\; \frac{1}{T}\sum_{k=-\infty}^{+\infty} X\!\big(j(\omega-k\omega_s)\big).\,} \tag{4}\]

In plain English: the spectrum of the sampled signal is just the original spectrum $X(j\omega)$ \textit{copied and pasted} at every integer multiple of $\omega_s$, and each copy is scaled by $1/T$.

\subsection*{Key Insight}

Sampling in time creates \textbf{periodic copies (replicas)} of the original spectrum $X(j\omega)$, shifted to every integer multiple of $\omega_s$ and scaled by $1/T$. The entire content of the sampling theorem comes down to one question: \textbf{do these copies overlap, or not?}

Let us visualize both cases. Suppose $X(j\omega)$ is a simple triangle living between $-\omega_M$ and $+\omega_M$.

\textit{[Figure: top panel labeled "$\omega_s > 2\omega_M$: No overlap" shows triangular spectrum replicas at $-\omega_s$, $0$, $\omega_s$ with a visible gap between them, height $1/T$. Bottom panel labeled "$\omega_s < 2\omega_M$: Overlap (aliasing)" shows the replicas intruding into each other near $\pm\omega_M$ with red "overlap!" annotations.]}

\section*{3 The Sampling Theorem}

\subsection*{The Sampling Theorem (Shannon / Nyquist)}

Let $x(t)$ be a \textbf{band-limited} signal with

\[X(j\omega) = 0 \quad \text{for}\quad |\omega| > \omega_M.\]

Then $x(t)$ is \textbf{uniquely determined} by its samples $x(nT)$, $n = 0, \pm 1, \pm 2, \ldots$, if

\[\boxed{\;\omega_s > 2\omega_M\;} \qquad \text{where}\quad \omega_s = \frac{2\pi}{T}.\]

Given these samples, $x(t)$ can be reconstructed exactly by passing an impulse train (whose impulse amplitudes equal the sample values) through an ideal lowpass filter with gain $T$ and cutoff frequency $\omega_c$ satisfying $\omega_M < \omega_c < \omega_s - \omega_M$.

The frequency $2\omega_M$ is called the \textbf{Nyquist rate}: the minimum sampling frequency you need. The frequency $\omega_M$ (half the Nyquist rate) is often called the \textbf{Nyquist frequency}.

\subsection*{Key Insight}

The sampling theorem says two things:

\begin{enumerate}
  \item \textbf{Uniqueness:} If you sample fast enough, no information is lost.
  \item \textbf{Reconstruction:} You can get $x(t)$ back exactly using an ideal lowpass filter.
\end{enumerate}

The condition $\omega_s > 2\omega_M$ is a \textit{strict} inequality. Sampling at exactly twice the highest frequency is \textit{not} sufficient in general (Example 22.3 shows why).

\subsection*{3.1 The Reconstruction System}

The ideal lowpass filter $H(j\omega)$ isolates the center copy of the spectrum and throws away all the other replicas.

\textit{[Figure: block diagram — $x(t)$ enters a multiplier driven by $p(t)$, producing $x_p(t)$, which feeds an LPF $H(j\omega)$ whose output is $x_r(t) = x(t)$.]}

With cutoff at $\omega_c = \omega_s/2$: $H(j\omega) = T$ for $|\omega| < \omega_s/2$, and $H(j\omega) = 0$ for $|\omega| > \omega_s/2$.

Here is the full picture of what happens at each stage in the frequency domain:

\textit{[Figure: four stacked spectra — (a) Original $X(j\omega)$: triangle of height 1 between $\pm\omega_M$; (b) After sampling $X_p(j\omega)$: replicas at $0,\pm\omega_s,\ldots$ each of height $1/T$; (c) Ideal LPF $H(j\omega)$: rectangle of height $T$ spanning the passband; (d) Recovered $X_r(j\omega) = X(j\omega)$ labeled "Perfect!".]}

\section*{4 Worked Examples}

\subsection*{Example 22.1(a): Nyquist Rate of a Sum of Sinusoids}

\textbf{Problem.} Find the Nyquist rate of $x(t) = 1 + \cos(2000\pi t) + \sin(4000\pi t)$.

The three frequency components are at $0$, $2000\pi$, and $4000\pi$ rad/s. The highest is $\omega_M = 4000\pi$ rad/s = 2000 Hz.

\[\text{Nyquist rate} = 2\omega_M = 8000\pi\ \text{rad/s} = 4000\ \text{Hz}.\]

We need to sample faster than 4000 times per second.

\textit{[Figure: magnitude spectrum $|X(j\omega)|$ showing DC impulse at 0, impulses at $\pm 2000\pi$, and impulses at $\pm 4000\pi$; the double-headed arrow "Nyquist rate = $2\omega_M$" spans from $-\omega_M$ to $+\omega_M$.]}

\subsection*{Example 22.1(b): Nyquist Rate of a Sinc Function}

\textbf{Problem.} Find the Nyquist rate of $x(t) = \dfrac{\sin(4000\pi t)}{\pi t}$.

We rewrite: $x(t) = 4000 \cdot \dfrac{\sin(4000\pi t)}{4000\pi t} = 4000\,\mathrm{sinc}(4000\,t)$, where $\mathrm{sinc}(u) \triangleq \sin(\pi u)/(\pi u)$.

Its Fourier transform is a \textbf{rectangle}: $X(j\omega) = 1$ for $|\omega| < 4000\pi$, and $0$ otherwise. So $\omega_M = 4000\pi$ and the Nyquist rate is $8000\pi$ rad/s = 4000 Hz.

\textit{[Figure: left — time domain sinc waveform; right — frequency domain rectangle spanning $-4000\pi$ to $4000\pi$.]}

\subsection*{Example 22.1(c): Nyquist Rate of a Squared Sinc}

\textbf{Problem.} Find the Nyquist rate of $x(t) = \left(\dfrac{\sin(4000\pi t)}{\pi t}\right)^{2}$.

First, recognize what is being squared. From part (b) we know that $\dfrac{\sin(4000\pi t)}{\pi t} = 4000\,\mathrm{sinc}(4000\,t)$ is a sinc function with bandwidth $\omega_M = 4000\pi$.

Now we are squaring that entire sinc function. Squaring in the time domain corresponds to convolving the spectrum with itself in the frequency domain (since $x^2(t) \leftrightarrow \frac{1}{2\pi}\,X(j\omega) * X(j\omega)$). The convolution of two rectangles, each of width $2 \times 4000\pi = 8000\pi$, produces a \textbf{triangle} of total width $2 \times 8000\pi = 16000\pi$. Therefore $\omega_M = 8000\pi$ and:

\[\text{Nyquist rate} = 2\omega_M = 16000\pi\ \text{rad/s} = 8000\ \text{Hz}.\]

The lesson: squaring a signal \textbf{doubles} the bandwidth!

\textit{[Figure: left — $X(j\omega)$ rectangle spanning $\pm\omega_M$; right — after squaring, $X * X$ triangle spanning $\pm 2\omega_M$.]}

\subsection*{Example 22.2: Checking Whether a Sampling Period Works}

\textbf{Problem.} A signal $x(t)$ is the output of an ideal lowpass filter with cutoff frequency $\omega_c = 1000\pi$ rad/s. Which of the following sampling periods guarantee that $x(t)$ can be perfectly reconstructed from its samples?

(a) $T = 0.5 \times 10^{-3}$ s
(b) $T = 2 \times 10^{-3}$ s
(c) $T = 10^{-4}$ s

\textbf{Step 1: Identify the highest frequency.}
Since $x(t)$ comes out of an ideal lowpass filter with cutoff $\omega_c = 1000\pi$ rad/s, we know $X(j\omega) = 0$ for $|\omega| > 1000\pi$. So the highest frequency present in $x(t)$ is:

\[\omega_M = 1000\pi\ \text{rad/s} \qquad (\text{equivalently, } f_M = 500\ \text{Hz}).\]

\textbf{Step 2: Compute the Nyquist rate.}
The sampling theorem requires $\omega_s > 2\omega_M$:

\[\omega_s > 2 \times 1000\pi = 2000\pi\ \text{rad/s} \qquad (\text{equivalently, } f_s > 1000\ \text{Hz}).\]

\textbf{Step 3: Convert to a constraint on $T$.}
Since $\omega_s = 2\pi/T$, the condition $\omega_s > 2000\pi$ becomes:

\[\frac{2\pi}{T} > 2000\pi \quad\Longrightarrow\quad T < \frac{2\pi}{2000\pi} = \frac{1}{1000} = 10^{-3}\ \text{s} = 1\ \text{ms}.\]

Any sampling period \textit{less than} 1 ms will work.

\textbf{Step 4: Check each option.}

\textit{[Figure: number line showing $T_\text{max} = 1$ ms as dashed boundary; green region "No aliasing" to the left with (c) 0.1 and (a) 0.5; red region "Aliasing" to the right with (b) 2.0.]}

\begin{itemize}
  \item \textbf{(a)} $T = 0.5$ ms $< 1$ ms: $\omega_s = 2\pi/(0.5 \times 10^{-3}) = 4000\pi > 2000\pi$. \checkmark This works. We are sampling at twice the Nyquist rate (oversampled by 2$\times$).
  \item \textbf{(b)} $T = 2$ ms $> 1$ ms: $\omega_s = 2\pi/(2 \times 10^{-3}) = 1000\pi < 2000\pi$. $\times$ This is too slow. The sampling frequency is only half of what we need, so aliasing will distort the signal.
  \item \textbf{(c)} $T = 0.1$ ms $< 1$ ms: $\omega_s = 2\pi/(10^{-4}) = 20000\pi \gg 2000\pi$. \checkmark This works easily. We are sampling at 10$\times$ the Nyquist rate.
\end{itemize}

\subsection*{Example 22.3: Sampling at Exactly the Nyquist Rate Fails}

\textbf{Problem.} Consider $x(t) = \sin\!\big(\tfrac{\omega_s}{2}t\big)$ sampled at frequency $\omega_s$ (exactly two samples per cycle).

\textbf{Solution.} The samples are $x(nT) = \sin\!\big(\tfrac{\omega_s}{2}\cdot\tfrac{2\pi n}{\omega_s}\big) = \sin(\pi n) = 0$ for every integer $n$. Every sample is zero! The sampler sees nothing.

\textit{[Figure: sinusoid with red sample markers all landing exactly on zero crossings, labeled "All samples = 0!"; sampling period $T$ indicated.]}

This is why $\omega_s$ must be \textit{strictly greater than} $2\omega_M$.

\section*{5 Reconstruction via Band-Limited Interpolation}

In the time domain, the output of the ideal lowpass reconstruction filter is:

\[x_r(t) \;=\; \sum_{n=-\infty}^{+\infty} x(nT)\,h(t - nT), \qquad h(t) = \frac{\sin(\pi t/T)}{\pi t/T}. \tag{5}\]

So the reconstruction formula is:

\[\boxed{\,x_r(t) \;=\; \sum_{n=-\infty}^{+\infty} x(nT)\,\frac{\sin\!\big[\pi(t-nT)/T\big]}{\pi(t-nT)/T}.\,} \tag{6}\]

Place a sinc at every sample point, scale it by the sample value, and add them all up. Each sinc equals 1 at its own sample time and 0 at every other sample time.

\textit{[Figure: several dashed sinc functions centered at $0, T, 2T, 3T, 4T$ weighted by their sample values, summing (solid black) to a smooth reconstructed waveform; legend lists "sum", "sinc at 0", "sinc at $T$", etc.]}

\subsection*{Key Insight}

Equation (6) is the \textbf{ideal band-limited interpolation formula}. In practice, the sinc extends to $\pm\infty$ (you would need infinitely many samples), so real systems use approximate methods: zero-order holds, linear interpolation, or windowed-sinc filters.

\section*{6 Practical Interpolation: Zero-Order and First-Order Holds}

\subsection*{6.1 Zero-Order Hold (Staircase)}

Hold each sample value constant until the next sample arrives. Uses a rectangular pulse $h_0(t)$ of width $T$:

\textit{[Figure: left — $h_0(t)$ rectangular pulse of height 1, width $T$; right — "Zero-order hold output": staircase approximation (red) overlaid on the original smooth curve (blue dashed).]}

\subsection*{6.2 First-Order Hold (Connect-the-Dots)}

Connect adjacent samples with straight lines. Uses a triangular pulse $h_1(t)$:

\textit{[Figure: left — $h_1(t)$ triangular pulse of height 1 with base from $-T$ to $T$; right — "First-order hold output": piecewise-linear reconstruction (solid) overlaid on the original smooth curve (dashed).]}

Neither method is exact. Both approximate the ideal sinc interpolation.

\section*{7 The Effect of Undersampling: Aliasing}

When $\omega_s < 2\omega_M$, the shifted replicas of $X(j\omega)$ overlap, and the original spectrum cannot be separated out by lowpass filtering. This distortion is called \textbf{aliasing}: high-frequency components masquerade as lower frequencies.

\subsection*{Warning}

Aliasing is \textbf{irreversible}. Once spectral replicas overlap, there is no way to undo the damage. The only way to prevent aliasing is to ensure $\omega_s > 2\omega_M$ \textit{before} sampling, or to pre-filter the signal to remove frequencies above $\omega_s/2$ (an \textbf{anti-aliasing filter}).

\subsection*{7.1 Aliasing of a Sinusoid}

\subsubsection*{Example 22.4: A 5 Hz Cosine Sampled at 6 Hz}

Consider $x(t) = \cos(2\pi \cdot 5\,t)$ sampled at $f_s = 6$ Hz. The Nyquist rate is $2f_0 = 10$ Hz. Since $f_s = 6 < 10$, we are undersampling. The aliased frequency is $f_\text{alias} = f_s - f_0 = 6 - 5 = 1$ Hz.

\textit{[Figure: 5 Hz cosine (blue) and 1 Hz cosine (red dashed) plotted on the same axis with red sample dots showing both curves passing through the same sample points; caption: "Both curves pass through the same sample points!"]}

The 5 Hz and 1 Hz signals produce the same samples. The lowpass filter picks the lower one, so the 5 Hz tone is aliased to 1 Hz.

\subsection*{7.2 Aliasing in the Frequency Domain}

The original spectrum has impulses at $\pm 5$ Hz. After sampling at 6 Hz, replicas appear at $\pm 5 + 6k$. The replicas at $\pm 1$ Hz land inside the lowpass filter:

\textit{[Figure: (a) Original spectrum with impulses at $\pm 5$ Hz. (b) After sampling at $f_s = 6$ Hz: impulses at $\pm 5$ plus additional circled impulses at $\pm 1$ Hz (labeled "Circled = aliased") inside the LPF passband box "$|f| < 3$".]}

\subsection*{7.3 Everyday Examples of Aliasing}

\textbf{Wagon-wheel effect in movies:} A movie camera captures frames at 24 fps. If a wagon wheel spins near 12 revolutions per second, it can appear to stand still or rotate backward. The rotation frequency exceeds half the frame rate.

\textbf{Moiré patterns:} Photographing a fine-striped shirt with a coarse pixel grid produces wavy interference patterns: spatial aliasing.

\section*{8 Discrete-Time Processing of Continuous-Time Signals}

One of the most powerful applications of sampling is using discrete-time systems (computers, DSP chips) to process continuous-time signals:

\textit{[Figure: block diagram — $x_c(t)$ into "C/D converter" (A-to-D) producing $x_d[n]$, into "Discrete-time system" producing $y_d[n]$, into "D/C converter" (D-to-A) producing $y_c(t)$.]}

The C/D converter creates $x_d[n] = x_c(nT)$. The D/C converter reconstructs $y_c(t)$ from $y_d[n]$ via lowpass filtering. If $x_c(t)$ is band-limited and the sampling theorem is satisfied, this entire cascade behaves like a continuous-time LTI system. The discrete-time frequency $\Omega$ and the continuous-time frequency $\omega$ are related by $\Omega = \omega T$, so the full range $|\omega| < \omega_s/2$ maps onto $|\Omega| < \pi$.

\section*{9 Summary of Key Formulas}

\begin{center}
\begin{tabular}{|l|l|}
\hline
\textbf{Concept} & \textbf{Formula / Condition} \\
\hline
Sampling period / frequency & $T = 2\pi/\omega_s$ \\
\hline
Nyquist rate & $2\omega_M$ (must be exceeded by $\omega_s$) \\
\hline
Sampling condition & $\omega_s > 2\omega_M$ (strict inequality!) \\
\hline
Spectrum of sampled signal & $X_p(j\omega) = \dfrac{1}{T}\sum_{k} X\!\big(j(\omega - k\omega_s)\big)$ \\
\hline
Sinc interpolation & $x_r(t) = \sum_{n} x(nT)\,\dfrac{\sin[\pi(t-nT)/T]}{\pi(t-nT)/T}$ \\
\hline
Aliased frequency & $f_\text{alias} = |f_s - f_0|$ (for $f_s/2 < f_0 < f_s$) \\
\hline
CT $\leftrightarrow$ DT frequency & $\Omega = \omega T$ \\
\hline
\end{tabular}
\end{center}

\section*{10 Summary}

\textit{[Figure: flow diagram — "Continuous-time signal $x(t)$" $\to$ "Sample at rate $\omega_s > 2\omega_M$" $\to$ "Discrete-time signal $x[n]$" $\to$ "Ideal LPF reconstruction" looping back to the continuous-time signal.]}

\textbf{Key takeaways:}

\begin{enumerate}
  \item \textbf{Sampling in time creates periodic copies in frequency.} The sampled spectrum is the original copy-pasted at every multiple of $\omega_s$, scaled by $1/T$.
  \item \textbf{The sampling theorem} guarantees perfect reconstruction when $\omega_s > 2\omega_M$. Reconstruction uses an ideal lowpass filter (sinc interpolation in time).
  \item \textbf{Aliasing} is the irreversible distortion from sampling too slowly. Higher frequencies fold back into lower ones. Use anti-aliasing filters \textit{before} sampling.
  \item \textbf{Practical reconstruction} uses zero-order holds, first-order holds, or other approximate filters.
  \item \textbf{Discrete-time processing of CT signals} works because sampling preserves all information in a band-limited signal. The link: $\Omega = \omega T$.
\end{enumerate}

\section*{11 Common Mistakes to Avoid}

\begin{enumerate}
  \item \textbf{Confusing Hz and rad/s.} $f_s > 2f_M$ (Hz) or $\omega_s > 2\omega_M$ (rad/s). Mixing units gives factor-of-$2\pi$ errors.
  \item \textbf{Assuming $\omega_s = 2\omega_M$ is sufficient.} It is not (Example 22.3).
  \item \textbf{Forgetting the $1/T$ scaling.} Replicas are scaled by $1/T$; the reconstruction filter compensates with gain $T$.
  \item \textbf{Thinking aliasing can be undone.} It cannot. Prevent it before sampling.
  \item \textbf{Claiming non-band-limited signals can't be sampled.} In practice, we band-limit with an anti-aliasing filter first.
\end{enumerate}

\section*{12 Practice Questions}

Test your understanding with these quick questions. Answers are at the bottom of the section.

\begin{enumerate}
  \item A signal $x(t)$ has $X(j\omega) = 0$ for $|\omega| > 5000\pi$. What is the Nyquist rate? What is the maximum allowable sampling period $T$?
  \item True or false: if $\omega_s = 2\omega_M$, the sampling theorem guarantees perfect reconstruction.
  \item In one sentence, explain \textit{why} sampling in time produces copies of the spectrum in the frequency domain.
  \item A signal is sampled and the resulting spectrum $X_p(j\omega)$ shows overlapping replicas. Can you undo the overlap with a clever filter? Why or why not?
  \item What is an anti-aliasing filter, and \textit{where} in the signal chain does it go (before or after the sampler)?
  \item Determine the Nyquist rate for $x(t) = \cos(600\pi t) + \cos(1800\pi t)$.
  \item A signal with $\omega_M = 3000\pi$ is sampled with $T = 2 \times 10^{-4}$ s. Compute $\omega_s$ and state whether aliasing occurs.
  \item You sample $x(t) = \cos(2\pi \cdot 900\,t)$ at $f_s = 1000$ Hz. What frequency does the reconstructed signal have?
  \item Repeat the previous question for $f_s = 1500$ Hz.
  \item A signal has Nyquist rate $\omega_0$. What is the Nyquist rate of $x(t-5)$? What about $x(2t)$?
  \item The ideal reconstruction filter has gain $T$ in its passband. Why is that factor of $T$ needed? (Hint: look at the $1/T$ scaling of the replicas.)
  \item A zero-order hold is used instead of an ideal lowpass filter for reconstruction. Is the result exact? What kind of error does it introduce?
  \item A movie camera shoots at 24 frames per second. A helicopter blade spins at 25 revolutions per second. Describe qualitatively what the blade looks like in the film.
  \item You are told that $X(j\omega) = 0$ for $|\omega| > 10{,}000\pi$. Someone samples $x(t)$ at $f_s = 8000$ Hz. Is the sampling theorem satisfied? If not, what is the minimum $f_s$ that works?
  \item Suppose $x(t)$ is \textit{not} band-limited (its spectrum extends to $\pm\infty$). Does this mean we can never sample it in practice? Explain what is done in real systems.
\end{enumerate}

\textit{Rogelio Gracia Otalvaro}

\section*{Practice Problems Summary}

\begin{itemize}
  \item \textbf{Example 22.1(a):} Compute the Nyquist rate of a sum of sinusoids $1 + \cos(2000\pi t) + \sin(4000\pi t)$ by identifying the highest frequency component.
  \item \textbf{Example 22.1(b):} Find the Nyquist rate of a sinc function $\sin(4000\pi t)/(\pi t)$ by recognizing its Fourier transform as a rectangle.
  \item \textbf{Example 22.1(c):} Determine the Nyquist rate of a squared sinc, using the fact that squaring in time doubles the bandwidth (rectangle convolved with itself becomes a triangle).
  \item \textbf{Example 22.2:} Decide which of three candidate sampling periods (0.5 ms, 2 ms, 0.1 ms) allow perfect reconstruction of a signal band-limited to $1000\pi$ rad/s.
  \item \textbf{Example 22.3:} Demonstrate that sampling $\sin(\omega_s t/2)$ at exactly $\omega_s$ yields all-zero samples, showing the strict inequality in the sampling theorem is required.
  \item \textbf{Example 22.4:} Show how a 5 Hz cosine sampled at 6 Hz aliases to a 1 Hz cosine because $f_s < 2f_0$.
  \item \textbf{Practice Question 1:} Compute Nyquist rate and maximum sampling period for a signal with $\omega_M = 5000\pi$.
  \item \textbf{Practice Question 2:} True/false check on whether $\omega_s = 2\omega_M$ guarantees perfect reconstruction.
  \item \textbf{Practice Question 3:} Explain in one sentence why sampling creates spectral copies.
  \item \textbf{Practice Question 4:} Explain whether overlapping replicas (aliasing) can be undone by filtering.
  \item \textbf{Practice Question 5:} Define an anti-aliasing filter and state its placement in the signal chain.
  \item \textbf{Practice Question 6:} Compute the Nyquist rate of $\cos(600\pi t) + \cos(1800\pi t)$.
  \item \textbf{Practice Question 7:} Given $\omega_M = 3000\pi$ and $T = 2 \times 10^{-4}$ s, compute $\omega_s$ and determine whether aliasing occurs.
  \item \textbf{Practice Question 8:} Identify the reconstructed frequency when $\cos(2\pi \cdot 900\,t)$ is sampled at 1000 Hz.
  \item \textbf{Practice Question 9:} Repeat the aliasing analysis with $f_s = 1500$ Hz.
  \item \textbf{Practice Question 10:} Determine the Nyquist rate of $x(t-5)$ and $x(2t)$ given that $x(t)$ has Nyquist rate $\omega_0$.
  \item \textbf{Practice Question 11:} Explain why the ideal reconstruction filter must have gain $T$ in its passband.
  \item \textbf{Practice Question 12:} Describe the error introduced when a zero-order hold replaces the ideal lowpass filter.
  \item \textbf{Practice Question 13:} Qualitatively describe the appearance of a 25 rev/s helicopter blade filmed at 24 fps.
  \item \textbf{Practice Question 14:} Determine whether $f_s = 8000$ Hz satisfies the sampling theorem for $\omega_M = 10{,}000\pi$, and find the minimum valid $f_s$.
  \item \textbf{Practice Question 15:} Explain how non-band-limited signals are handled in practice.
\end{itemize}

\end{document}
